\begin{tomtat}
\begin{itemize}
  \item Completeness, sensitivity và \tn{tính tuyến tính}{linearity} loại một số hành vi của
  attribution, nhưng khi đi cùng nhau chúng có thể bó lớp phương pháp vào dạng
  gần Gradient$\mathbin{\times}$Input.
  \item Khung của chương bắt đầu bằng attribution cho \tn{hàm chỉ báo}{indicator function}, rồi mở rộng
  sang mô hình liên tục bằng tính tuyến tính và một điều kiện ổn định theo mô
  hình.
  \item Một attribution trong khung có dạng tích phân của $f$ theo độ đo
  $\mu_{j,x}$. Độ đo ghi lựa chọn về những đầu vào nào được lấy trung bình.
  \item Nhiều phương pháp dựa trên feature removal và partial dependence plot
  xuất hiện như các lựa chọn độ đo khác nhau. Khung thống nhất phép viết, không
  chọn thay giả định miền ứng dụng.
\end{itemize}
\end{tomtat}

\cauhoi
\begin{cauhoids}
  \item Vì sao hai attribution cùng thỏa tính đầy đủ vẫn có thể trả lời hai
  câu hỏi khác nhau về một đầu vào?
  \item Trong phương trình~\ref{eq:ch06-attribution-do-do}, hãy mô tả bằng lời
  một độ đo giữ $x_j$ cố định và lấy trung bình các tọa độ còn lại. Giả định
  nào về dữ liệu xuất hiện trong mô tả đó?
  \item Một mạng ReLU là affine trên từng vùng kích hoạt cố định. Vì sao nhận
  xét này hữu ích cho công thức attribution, nhưng chưa giải quyết việc chọn
  độ đo?
  \item Đọc \enquote{Feature Attribution from First Principles} của
  Taimeskhanov và Garreau~\autocite{p17firstprinciples}. Tác giả cần những điều
  kiện nào để đi từ attribution của hàm chỉ báo đến tích phân cho mô hình liên
  tục?
\end{cauhoids}
